\chapter{ГИБРИДНЫЙ МЕТОД AMPLE-AMR}\label{sec:ch5}

Во второй главе была построена доменная и математическая модель системы
автономных мобильных роботов с поддержкой периферийных вычислений. В третьей
главе был разработан VCG-подобный аукционный слой распределения задач при
фиксированных операционных режимах периферийных узлов. В четвёртой главе
задача выбора этих режимов была формализована как децентрализованный
частично наблюдаемый марковский процесс принятия решений и обосновано
применение алгоритма QMIX.

Настоящая глава объединяет эти компоненты в единый гибридный метод
\textbf{AMPLE-AMR} (\textit{Adaptive Multi-agent Policy Learning with Edge Auctions for
Autonomous Mobile Robots}). Метод предназначен для адаптивного распределения вычислительных
задач автономных мобильных роботов между периферийными вычислительными
узлами. Его ключевая идея состоит в разделении функций между двумя слоями:
\begin{enumerate}
	\item обучаемый многоагентный слой выбирает операционные режимы
	периферийных узлов;
	\item VCG-подобный аукционный слой распределяет задачи между узлами в
	рамках выбранных ресурсных ограничений.
\end{enumerate}

Такое разделение позволяет избежать методологической ошибки, при которой
обучаемая политика начинает управлять ставками или оценками ценности задач.
В предлагаемом методе ставки и оценки ценности не обучаются. Обучается только
политика выбора режимов периферийных узлов: допуска, резервирования,
приоритизации и защиты от перегрузки.

В главе также вводится масштабируемая версия метода \textbf{C-AMPLE-AMR} (\textit{Clustered Adaptive Multi-agent Policy Learning with Edge Auctions
for Autonomous Mobile Robots}). Она предназначена для случаев, когда глобальный аукционный раунд становится
слишком дорогим по времени вычислений или по объёму коммуникации. В этой
версии граф периферийной инфраструктуры разбивается на локальные кластеры,
внутри которых проводятся локальные аукционы.

\section{Концепция метода}

\subsection{Общие принципы}

Метод AMPLE-AMR строится на следующих принципах.

Первый принцип --- разделение адаптации и распределения. Адаптация к
динамике среды выполняется обучаемой многоагентной политикой, а собственно
распределение задач выполняется аукционным механизмом. Это позволяет
использовать обучение для управления режимами узлов, не разрушая смысл
VCG-подобного механизма как правила распределения при фиксированных
ограничениях \cite{Vickrey1961,Clarke1971,Groves1973}.

Второй принцип --- локальность исполнения. На этапе исполнения каждый
периферийный узел принимает решение о своём режиме на основе локального
наблюдения и собственной истории. Это соответствует парадигме
централизованного обучения и децентрализованного исполнения, используемой в
кооперативном многоагентном обучении \cite{Rashid2018QMIX}.

Третий принцип --- экономическая интерпретируемость распределения. Аукционный
слой формирует не только назначение задач, но и показатели внешнего эффекта,
платежей, дефицита ресурсов и вклада в общественное благосостояние. Эти
величины используются как признаки и компоненты вознаграждения, но не как
обучаемые ставки.

Четвёртый принцип --- масштабируемость через локализацию распределения. При
росте числа роботов, задач и периферийных узлов один глобальный аукцион может
стать узким местом. Поэтому в методе C-AMPLE-AMR вводится разбиение графа
инфраструктуры на кластеры и проведение локальных аукционов внутри них.

Пятый принцип --- явное указание границ гарантий. VCG-подобные свойства
стратегической совместимости и эффективности рассматриваются только для
отдельного аукционного раунда при выполнении стандартных предпосылок. Они не
переносятся автоматически на всю повторяющуюся систему с обучением.

\subsection{Схема метода}

Пусть в момент времени $t$ имеется множество задач $\mathcal{T}(t)$,
порожденных роботами, и множество периферийных узлов $\mathcal{N}$. Один шаг
метода AMPLE-AMR состоит из следующих этапов:
\begin{enumerate}
	\item сбор локальных наблюдений периферийных узлов;
	\item выбор операционных режимов узлов обучаемой многоагентной политикой;
	\item вычисление экспонируемой ёмкости каждого узла;
	\item формирование заявок роботов и параметров узлов для аукционного
	слоя;
	\item решение задачи выбора победителей;
	\item вычисление VCG-подобных платежей и внутренних показателей внешнего
	эффекта;
	\item исполнение назначений, обновление очередей и загрузки;
	\item формирование командного вознаграждения;
	\item сохранение перехода для последующего обучения.
\end{enumerate}

Формально обучаемая политика задаёт вектор режимов
\[
\kappa_t =
\left(
\kappa_1(t),\ldots,\kappa_k(t)
\right),
\]
где $\kappa_j(t)$ --- режим периферийного узла $n_j$.

После выбора режима определяется экспонируемая ёмкость:
\[
\bar{A}^{r}_j(t,\kappa_j)
=
\eta^{r}_j(\kappa_j)
\left(
C^{r}_j - L^{r}_j(t)
\right),
\quad
r \in \{cpu,mem,gpu\}.
\]

Затем аукционный слой решает задачу:
\[
x^*(t)
\in
\arg\max_{x \in \mathcal{X}(t,\kappa(t))}
SW(x,t),
\]
где $\mathcal{X}(t,\kappa(t))$ --- множество допустимых назначений с учётом
достижимости, отказов, сроков обработки и экспонируемой ёмкости узлов.

Именно такая последовательность является основной архитектурной идеей метода:
обучаемая политика меняет не ставки и не ценности задач, а допустимые ресурсные
условия, в которых проводится аукцион.

\subsection{Роли участников и поток данных}

В методе AMPLE-AMR участвуют три логических уровня.

\textbf{Первый уровень} --- автономные мобильные роботы. Они формируют вычислительные
задачи, указывают класс задачи, приоритет, допустимый срок обработки, объём
данных и ресурсные требования. На основе этих параметров вычисляется оценка
ценности задачи. Роботы не являются обучаемыми агентами в MARL-контуре.

\textbf{Второй уровень} --- периферийные вычислительные узлы. Каждый узел является
обучаемым агентом, который выбирает операционный режим. Узел не назначает
задачи напрямую, а определяет, какую часть ресурсов предоставить для текущего
аукционного распределения.

\textbf{Третий уровень} --- аукционный слой. Он получает задачи, оценки ценности,
состояние достижимости, стоимости обслуживания и экспонируемую ёмкость
узлов. На этой основе он вычисляет распределение задач и VCG-подобные
платежи.

\begin{figure}[h]
	\centering
	\includegraphics[width=10cm, height=10cm]{ma-vcg-qmix}
	\caption{Архитектура AMPLE-AMR}
	\label{img:ma-vcg-qmix}
\end{figure}

Поток данных между уровнями можно представить следующим образом:
\[
\text{AMR tasks}
\rightarrow
\text{node modes}
\rightarrow
\text{auction allocation}
\rightarrow
\text{execution}
\rightarrow
\text{learning experience}.
\]

В более детальной форме:
\begin{enumerate}
	\item роботы передают параметры задач;
	\item периферийные узлы формируют локальные наблюдения;
	\item локальные политики узлов выбирают режимы;
	\item режимы определяют экспонируемые ресурсы;
	\item аукционный слой распределяет задачи;
	\item результаты аукциона формируют назначения и экономические признаки;
	\item среда обновляет состояние;
	\item обучаемый контур получает переход и командное вознаграждение.
\end{enumerate}

\subsection{Аукционные признаки}

Аукционный слой для каждого узла $n_j$ формирует вектор признаков, которые используются обучаемым контуром в последующих шагах:
\[
g^j_t =
\left(
n^{alloc}_j(t),
\overline{SW}_j(t),
\overline{p}_j(t),
q^N_j(t),
\sigma^{cpu}_j(t),
\sigma^{mem}_j(t),
\sigma^{gpu}_j(t),
d^{drop}_j(t)
\right),
\]
где:
\begin{enumerate}
	\item $n^{alloc}_j(t)$ --- число задач, назначенных периферийному узлу;
	\item  $\overline{SW}_j(t)$ --- суммарный и средний вклад назначенных задач в общественное благосостояние;
	\item $\overline{p}_j(t)$ --- VCG-подобные платежи роботов, чьи задачи были назначены узлу;
	\item  $q^N_j(t)$ --- внутренняя компенсация или теневая ценность узла;
	\item $\sigma^r_j(t)$ --- доля использованной экспонируемой ёмкости по ресурсу $r$;
	\item $d^{drop}_j(t)$ --- доля отклонённых или отложенных задач.
\end{enumerate}

Эти признаки являются экономически интерпретируемыми сигналами. Они
отражают дефицитность ресурсов, внешний эффект и вклад узла в текущее
распределение. 

\subsection{Функция вознаграждения}

Метод AMPLE-AMR использует командное вознаграждение, одинаковое для всех
периферийных узлов. Это соответствует кооперативной постановке: узлы должны
оптимизировать качество всей AMR-системы, а не только свои локальные
показатели.

\textbf{Базовая часть} вознаграждения $r^{base}_t$ определяется через функцию~\eqref{eq:ch4_reward_base}, куда входит нормированное общественное благосостояние и штрафы.

\textbf{Аукционная часть} $r^{auc}_t$ вознаграждения задаётся через фукнцию~\eqref{eq:ch4_reward_auc} как взвешенная сумма нормированных внешних и внутренних показателей полученного распределения.

Итоговое вознаграждение:
\begin{equation}
	\tilde{r}_t =
	(1-\lambda_{auc}) r_t
	+
	\lambda_{auc} r^{auc}_t,
	\quad
	\lambda_{auc}\in[0,1].
	\label{eq:ch5_reward}
\end{equation}

Параметр $\lambda_{auc}$ управляет силой влияния аукционных признаков на
обучение. При $\lambda_{auc}=0$ метод вырождается в вариант многоагентного
управления без экономического усиления. При слишком большом значении
$\lambda_{auc}$ возникает риск переориентации политики на платежные показатели,
поэтому в экспериментальной части требуется анализ чувствительности к этому
параметру.

\section{Режимы функционирования метода}

Алгоритм AMPLE-AMR функционирует в двух режимах: режиме онлайн-исполнения (Online Phase) и режиме оффлайн-обучения (Offline Phase).

\subsection{Режим 1: Онлайн-исполнение}

В режиме онлайн-исполнения используется уже обученная политика. Один шаг
метода выполняется следующим образом:

\begin{enumerate}
	\item \textbf{Поступление задач.} Поступают вычислительные задачи от автономных мобильных роботов.
	\item \textbf{Сбор наблюдений.} Каждый периферийный узел $n_j$ получает локальное наблюдение $o^j_t$.
	\item \textbf{Выбор локальных режимом.} Локальная сеть агента выбирает режим:
	\[
	\kappa_j(t)
	=
	\arg\max_{\kappa_j \in \mathcal{K}_j}
	Q_j(h^j_t,\kappa_j).
	\]
	\item \textbf{Определение ёмкости узлов.} На основе выбранного режима вычисляется экспонируемая ёмкость
	$\bar{A}_j(t,\kappa_j)$.
	\item \textbf{Запуск аукционного раунда.} Аукционный слой решает задачу выбора победителей и формирует распределение $x^*(t)$.
	\item \textbf{Исполнение.} Назначенные задачи передаются в очереди соответствующих узлов или запускаются на исполнение.
	\item \textbf{Наблюдение результата.} Среда обновляет состояние: загрузку, очереди, задержки, завершённые задачи и отказовые признаки.
	\item \textbf{Сохранение перехода.} Формируются показатели текущего шага и сохраняются журналы для последующего обучения или анализа.
\end{enumerate}

Онлайн-исполнение не требует глобального пересчёта политики. Глобальная
информация используется аукционным слоем для текущего распределения и
серверным контуром для анализа, но локальный выбор режима осуществляется
каждым узлом самостоятельно.

\subsection{Режим 2: Оффлайн-обучение}

Оффлайн-обучение проводится в серверном контуре, облаке, цифровом двойнике
или имитационной среде. В отличие от онлайн-исполнения, здесь допускается
использование глобального состояния системы, совместных действий всех агентов,
полных результатов аукционных раундов и командного вознаграждения.

Обучающий цикл включает:
\begin{enumerate}
	\item генерацию или воспроизведение сценария работы AMR-системы;
	\item выполнение последовательности шагов AMPLE-AMR;
	\item накопление переходов в буфере опыта;
	\item выбор мини-пакетов переходов;
	\item расчёт временно-разностных целей;
	\item обновление локальных сетей агентов и сети смешивания QMIX;
	\item периодическое обновление целевых сетей;
	\item проверку качества политики на отложенных сценариях.
\end{enumerate}

Переход, сохраняемый в буфере опыта, имеет вид:
\[
e_t =
\left(
s_t,
\{o^j_t\}_{j=1}^{k},
\kappa_t,
r_t,
s_{t+1},
\{o^j_{t+1}\}_{j=1}^{k},
done_t
\right).
\]

Сеть QMIX обучается минимизацией функции потерь:
\[
\mathcal{L}(\theta)
=
\mathbb{E}_{e_t\sim\mathcal{B}}
\left[
\left(
y_t -
Q_{tot}(s_t,\kappa_t;\theta)
\right)^2
\right],
\]
где
\[
y_t =
r_t
+
\gamma(1-done_t)
Q^{-}_{tot}(s_{t+1},\kappa'_{t+1}).
\]

После завершения обучения параметры локальных сетей распространяются на
периферийные узлы. Повторное обучение может запускаться периодически, при
существенном изменении нагрузки или после обновления конфигурации
инфраструктуры.

\section{Псевдокод метода AMPLE-AMR}

Ниже приведена алгоритмическая схема метода.

\textbf{Алгоритм 5.1. Метод AMPLE-AMR}

\begin{enumerate}
	\item Инициализировать локальные сети агентов, сеть смешивания QMIX,
	целевые сети и буфер опыта.
	\item Для каждого эпизода обучения выполнить последовательность шагов.
	\item На шаге $t$ получить множество задач $\mathcal{T}(t)$ от роботов.
	\item Для каждого периферийного узла $n_j$ сформировать локальное
	наблюдение $o^j_t$.
	\item Для каждого узла выбрать режим $\kappa_j(t)$ с помощью локальной
	политики.
	\item Вычислить экспонируемую ёмкость $\bar{A}_j(t,\kappa_j)$.
	\item Сформировать вход аукционного слоя: задачи, оценки ценности,
	стоимости обслуживания, достижимость и экспонируемую ёмкость.
	\item Решить задачу выбора победителей и получить распределение $x^*(t)$.
	\item Вычислить VCG-подобные платежи и агрегированные аукционные признаки.
	\item Исполнить назначение задач в среде.
	\item Обновить очереди, загрузку, сетевые задержки и состояние отказов.
	\item Рассчитать командное вознаграждение $r_t$.
	\item Сохранить переход $e_t$ в буфер опыта.
	\item Если выполняется условие обновления, выполнить шаг обучения QMIX.
	\item После завершения обучения использовать локальные политики для
	онлайн-выбора режимов узлов.
\end{enumerate}

\section{Редукции метода к базовым вариантам}

Метод AMPLE-AMR содержит несколько базовых вариантов, которые используются
для последующей экспериментальной проверки вклада отдельных компонентов.
Важно подчеркнуть, что в предлагаемой архитектуре обучаемая многоагентная
политика не выполняет непосредственное назначение задач на узлы. Она выбирает
операционные режимы периферийных узлов, а само распределение задач выполняет
отдельный слой назначения. Поэтому при исключении аукционного слоя необходимо
явно задать альтернативное правило распределения.

Если отключить обучаемый контур и зафиксировать режимы периферийных узлов,
метод вырождается в повторяющееся аукционное распределение:
\[
\text{AMPLE-AMR}
\rightarrow
\text{Repeated VCG-like allocation}.
\]
В этом варианте периферийные узлы всегда используют заранее заданные
режимы, например нормальный или консервативный, а распределение задач в
каждом раунде выполняется VCG-подобным механизмом.

Если сохранить обучаемый выбор режимов, но заменить VCG-подобный аукционный
слой эвристическим правилом назначения, получается вариант:
\[
\text{AMPLE-AMR}
\rightarrow
\text{QMIX + heuristic allocation}.
\]
В этом случае многоагентная политика по-прежнему выбирает режимы допуска,
резервирования и обслуживания периферийных узлов, но задачи назначаются не
через максимизацию общественного благосостояния и VCG-платежи, а с помощью
заданной эвристики. Такой эвристикой может быть выбор наименее загруженного
доступного узла, выбор узла с минимальной оценкой задержки, жадное назначение
по максимальному чистому вкладу или комбинированное правило, учитывающее
задержку, загрузку и дедлайн задачи.

Если сохранить и обучаемый выбор режимов, и VCG-подобный аукционный слой,
получается полный метод:
\[
\text{AMPLE-AMR}
=
\text{QMIX-based node mode control}
+
\text{VCG-like task allocation}.
\]
Именно этот вариант является основным в диссертации: обучаемая политика
адаптирует ресурсные режимы периферийных узлов к динамике среды, а
аукционный слой выполняет экономически интерпретируемое распределение задач
в рамках выбранных ограничений.

Наконец, если в полном методе дополнительно включить кластеризацию графа
периферийной инфраструктуры и проводить локальные аукционы внутри кластеров,
получается масштабируемый вариант:
\[
\text{C-AMPLE-AMR} = \text{clustered AMPLE-AMR}.
\]

Таким образом, для экспериментального сравнения целесообразно использовать
следующие варианты:
\begin{enumerate}
	\item фиксированные режимы узлов и VCG-подобное распределение;
	\item фиксированные режимы узлов и эвристическое распределение;
	\item обучаемые режимы узлов и эвристическое распределение;
	\item обучаемые режимы узлов и VCG-подобное распределение, то есть
	полный AMPLE-AMR;
	\item кластерный вариант C-AMPLE-AMR.
\end{enumerate}

\section{Проблема масштабируемости глобального аукциона}

Глобальный вариант AMPLE-AMR предполагает, что аукционный слой распределяет
задачи по всей системе. При малом и среднем числе роботов такой подход
приемлем. Однако при увеличении масштаба системы растут:
\begin{enumerate}
	\item число заявок;
	\item число допустимых пар задача--узел;
	\item размер задачи выбора победителей;
	\item число повторных решений для вычисления платежей;
	\item объём коммуникации для сбора состояния;
	\item задержка между поступлением задачи и принятием решения.
\end{enumerate}

Обозначим число задач текущего раунда через $d_t$, число узлов через $k$, а
число роботов через $m$. Пусть $T^{WDP}(d_t,k)$ --- время решения задачи выбора
победителей. Тогда прямое вычисление VCG-подобных платежей имеет порядок
\[
T^{VCG}(t)
=
\mathcal{O}
\left(
(m+k+1)T^{WDP}(d_t,k)
\right).
\]

К этой величине добавляется коммуникационная составляющая:
\[
T^{comm}(t)
=
T^{uplink}(m,d_t)
+
T^{state}(k)
+
T^{downlink}(m,k),
\]
где учитываются передача заявок, сбор состояния узлов и рассылка результатов
назначения.

Глобальный режим становится проблемным, если
\[
T^{VCG}(t)+T^{comm}(t)
>
\rho \Delta^{slot},
\]
где $\Delta^{slot}$ --- допустимая длительность управляющего слота, а
$\rho \in (0,1)$ --- максимально допустимая доля времени, отводимая на
служебное распределение.

Это условие используется как один из критериев перехода к кластерному
варианту C-AMPLE-AMR.

\section{Кластерная версия C-AMPLE-AMR}

Метод C-AMPLE-AMR предназначен для повышения масштабируемости за счёт
локализации аукционного распределения. Основная идея состоит в том, чтобы
разбить граф инфраструктуры на кластеры и проводить распределение задач внутри
каждого кластера отдельно.

Пусть граф инфраструктуры задан как~\eqref{eq:ch2_graph}
\[
G(t)=(\mathcal{V},\mathcal{E},w(t)),
\quad
\mathcal{V}=\mathcal{R}\cup\mathcal{N}.
\]

Для целей кластеризации вводится \textbf{аффинный вес}:
\begin{equation}
	\alpha_{uv}(t) = \frac{\nu_{uv}(t)}{\varepsilon + w_{uv}(t)},
	\label{eq:ch5_affine_value}
\end{equation}
где $w_{uv}(t)$ --- стоимость или задержка связи, $\nu_{uv}(t)$ --- частота или важность взаимодействия между вершинами, а $\varepsilon>0$ --- малый стабилизирующий коэффициент (предотвращает деление на ноль). 

Большое значение $\alpha_{uv}(t)$ означает, что вершины $u$ и $v$ сильно
связаны и предпочтительно должны попадать в один кластер. Это важно, поскольку
во многих алгоритмах выделения сообществ веса рёбер трактуются как сила связи,
а не как стоимость взаимодействия. Поэтому исходные задержки или стоимости
нельзя напрямую использовать как аффинные веса без преобразования.

Разбиение графа имеет вид:
\begin{equation}
	\mathcal{P}(t) = \{P_1(t),P_2(t),\ldots,P_s(t)\},
	\label{eq:ch5_clusters}
\end{equation}
где каждый кластер $P_\ell(t)$ содержит подмножество роботов и периферийных
узлов:
\begin{equation}
	P_\ell(t)
	=
	\mathcal{R}_\ell(t)
	\cup
	\mathcal{N}_\ell(t).
	\label{eq:ch5_cluster_nodes}
\end{equation}

\textbf{Цель кластеризации} состоит в том, чтобы:
\begin{enumerate}
	\item сохранить сильносвязанные пары робот--узел внутри одного кластера;
	\item уменьшить число межкластерных назначений;
	\item сбалансировать нагрузку между кластерами;
	\item ограничить размер локальных аукционов;
	\item снизить коммуникационные и вычислительные накладные расходы.
\end{enumerate}

Для построения разбиения могут использоваться методы разбиения графов и
выделения сообществ, включая METIS-подобные алгоритмы и методы модульности
\cite{Karypis1998Metis,Blondel2008Louvain}. Выбор конкретного алгоритма
зависит от требований реализации: необходимости баланса размеров кластеров,
скорости обновления разбиения и устойчивости к малым изменениям графа.

\subsection{Локальные аукционы внутри кластеров}

После получения разбиения $\mathcal{P}(t)$ глобальная задача распределения
заменяется набором локальных задач. Для кластера $P_\ell$ обозначим множество
роботов через $\mathcal{R}_\ell$, множество узлов через $\mathcal{N}_\ell$, а
множество задач через
\begin{equation}
	\mathcal{T}_\ell(t)
	=
	\bigcup_{r_i\in\mathcal{R}_\ell}
	\mathcal{T}_i(t).
	\label{eq:ch5_cluster_tasks}
\end{equation}

Локальный аукцион решает:
\begin{equation}
	x^*_\ell(t)
	\in
	\arg\max_{x_\ell \in \mathcal{X}_\ell(t,\kappa(t))}
	SW_\ell(x_\ell,t),
	\label{eq:ch5_cluster_alloc}
\end{equation}
где
\begin{equation}
	SW_\ell(x_\ell,t)
	=
	\sum_{\tau\in\mathcal{T}_\ell(t)}
	\sum_{n_j\in\mathcal{N}_\ell}
	x_{\tau j}(t)
	\left[
	U_{i(\tau)}(\tau,n_j,t)
	-
	C_j(\tau,r_{i(\tau)},t)
	\right].
	\label{eq:ch5_cluster_sw}
\end{equation}

Итоговое распределение в кластерном режиме:
\begin{equation}
	x^{C}(t)
	=
	\bigcup_{\ell=1}^{s}
	x^*_\ell(t).
	\label{eq:ch5_final_alloc}
\end{equation}

Такое распределение в общем случае может уступать глобальному оптимуму по
общественному благосостоянию, поскольку часть межкластерных назначений
исключается или ограничивается. Поэтому C-AMPLE-AMR должен использоваться
тогда, когда выигрыш по накладным расходам превышает возможную потерю
качества распределения.

\subsection{Межкластерные назначения}

В базовом варианте C-AMPLE-AMR задачи распределяются только внутри своих
кластеров. Однако в некоторых случаях это может приводить к потере качества:
например, если в кластере нет свободного узла, а соседний кластер имеет
избыточную ёмкость.

Для таких ситуаций можно ввести ограниченный механизм межкластерного
резерва. Пусть $\mathcal{B}_\ell(t)$ --- множество соседних кластеров,
доступных для кластера $P_\ell$. Тогда задача $\tau\in\mathcal{T}_\ell(t)$
может быть назначена не только на узлы $\mathcal{N}_\ell$, но и на ограниченное
число узлов из соседних кластеров:
\[
\mathcal{N}^{cand}_\ell(t)
=
\mathcal{N}_\ell(t)
\cup
\bigcup_{P_b\in\mathcal{B}_\ell(t)}
\mathcal{N}^{border}_b(t),
\]
где $\mathcal{N}^{border}_b(t)$ --- подмножество пограничных или резервных
узлов соседнего кластера.

Межкластерное назначение штрафуется дополнительной сетевой стоимостью:
\[
C^{cross}_{\tau j}(t)
=
C_j(\tau,r_{i(\tau)},t)
+
\rho^{cross} d_{\ell b}(t),
\]
где $d_{\ell b}(t)$ --- стоимость взаимодействия между кластерами, а
$\rho^{cross}$ --- коэффициент штрафа.

Такой механизм позволяет сохранить локальность распределения, но не блокирует
редкие выгодные межкластерные назначения.

\subsection{Критерий перехода к кластерному режиму}

Кластерный режим не должен использоваться всегда. Если система мала, один
глобальный аукцион может давать лучшее качество распределения без существенных
накладных расходов. Поэтому в методе вводится критерий перехода.

Пусть
\[
H^{global}(t)
=
T^{auction}_{global}(t)
+
T^{comm}_{global}(t)
\]
--- накладные расходы глобального режима.

Пусть
\[
H^{cluster}(t)
=
\max_{\ell}
\left(
T^{auction}_{\ell}(t)
+
T^{comm}_{\ell}(t)
\right)
+
T^{partition}(t)
+
T^{sync}(t)
\]
--- накладные расходы кластерного режима, где учитываются локальные аукционы,
разбиение графа и межкластерная синхронизация.

Переход к кластерному режиму целесообразен, если выполняются два условия:
\[
H^{global}(t)
>
\rho \Delta^{slot},
\]
и
\[
H^{cluster}(t)
+
\lambda_{loss} \widehat{L}^{welfare}(t)
<
H^{global}(t),
\]
где $\widehat{L}^{welfare}(t)$ --- оценка потери общественного благосостояния
из-за локализации распределения, а $\lambda_{loss}$ --- коэффициент приведения
потерь качества к шкале накладных расходов.

На практике $\widehat{L}^{welfare}(t)$ может оцениваться по малому числу
контрольных глобальных аукционов, по исторической статистике или по величине
межкластерного среза графа:
\[
Cut(\mathcal{P}(t))
=
\sum_{\ell\neq b}
\sum_{u\in P_\ell,\,v\in P_b}
\alpha_{uv}(t).
\]

Чем больше межкластерный срез, тем выше риск потери качества при локальных
аукционах. Поэтому кластеризация должна не только уменьшать размер локальных
задач, но и сохранять сильные связи внутри кластеров.

\subsection{Псевдокод C-AMPLE-AMR}

\textbf{Алгоритм 5.2. Кластерная версия C-AMPLE-AMR}

\begin{enumerate}
	\item Получить текущий граф инфраструктуры $G(t)$.
	\item Вычислить аффинные веса $\alpha_{uv}(t)$.
	\item Оценить накладные расходы глобального режима $H^{global}(t)$.
	\item Если глобальный режим допустим по времени, выполнить AMPLE-AMR без
	кластеризации.
	\item Иначе построить разбиение $\mathcal{P}(t)$.
	\item Для каждого периферийного узла выбрать операционный режим с помощью
	локальной политики.
	\item В каждом кластере вычислить экспонируемую ёмкость узлов.
	\item Провести локальный VCG-подобный аукцион внутри каждого кластера.
	\item При необходимости обработать ограниченное число межкластерных
	назначений через пограничные узлы.
	\item Исполнить полученные назначения.
	\item Обновить состояние среды, аукционные признаки и командное
	вознаграждение.
	\item Сохранить переход для последующего обучения.
\end{enumerate}

\section{Анализ вычислительной сложности}

В глобальном варианте AMPLE-AMR стоимость аукционного слоя определяется
числом задач $d_t$ и числом узлов $k$:
\[
T^{global}_{auction}(t)
=
T^{WDP}(d_t,k).
\]
С учётом вычисления платежей:
\[
T^{global}_{VCG}(t)
=
\mathcal{O}
\left(
(m+k+1)T^{WDP}(d_t,k)
\right).
\]

В кластерном варианте пусть система разбита на $s$ кластеров, а в кластере
$P_\ell$ имеется $d_\ell$ задач и $k_\ell$ узлов. Тогда стоимость локальных
аукционов:
\[
T^{cluster}_{auction}(t)
=
\max_{\ell=1,\ldots,s}
T^{WDP}(d_\ell,k_\ell)
\]
при параллельном выполнении локальных аукционов или
\[
T^{cluster}_{auction}(t)
=
\sum_{\ell=1}^{s}
T^{WDP}(d_\ell,k_\ell)
\]
при последовательном выполнении.

Если кластеры сбалансированы, то приблизительно
\[
d_\ell \approx \frac{d_t}{s},
\quad
k_\ell \approx \frac{k}{s}.
\]
Для комбинаторной задачи выбора победителей это может дать существенное
снижение времени расчёта, поскольку $T^{WDP}$ обычно растёт сверхлинейно по
размеру задачи.

Стоимость обучаемого контура для выбора режимов имеет порядок
\[
T^{QMIX}_{infer}(t)
=
\mathcal{O}(k|\mathcal{K}|),
\]
если размер локального пространства режимов равен $|\mathcal{K}|$. При
фиксированном числе режимов этот компонент растёт линейно по числу узлов и
обычно не является основным узким местом.

Следовательно, главный выигрыш C-AMPLE-AMR связан именно с уменьшением
размерности аукционного слоя, а не с ускорением локального вывода
многоагентной политики.

\section{Свойства метода AMPLE-AMR}

Метод AMPLE-AMR обладает следующими свойствами.

Во-первых, он сохраняет интерпретируемость распределения. Назначение задач
строится через явную задачу максимизации общественного благосостояния, а
платежи имеют смысл внешнего эффекта.

Во-вторых, он допускает адаптацию к динамике среды. Периферийные узлы
обучаются выбирать режимы, учитывая очереди, загрузку, задержки, результаты
предыдущих аукционов и штрафы за нарушения сроков.

В-третьих, он поддерживает децентрализованное исполнение обучаемой политики.
Каждый узел выбирает свой режим по локальному наблюдению, а централизованное
обучение может выполняться отдельно в серверном контуре или цифровом
двойнике.

В-четвёртых, метод не обучает стратегию заявок. Это принципиально отделяет
его от подходов, где обучение используется для изменения ставок участников.

В-пятых, метод допускает масштабируемое расширение. C-AMPLE-AMR позволяет
заменить один глобальный аукцион набором локальных аукционов в кластерах.

В-шестых, метод имеет контролируемые частные случаи: только аукцион, только
многоагентное управление, эвристическое распределение и полный гибридный
вариант. Это удобно для экспериментальной проверки вклада каждого компонента.

\section{Ограничения метода}

Несмотря на преимущества, метод AMPLE-AMR имеет ряд ограничений.

Во-первых, VCG-подобные свойства относятся только к локальному аукционному
раунду при фиксированных режимах узлов. Вся повторяющаяся система с
обучением не является классическим VCG-механизмом.

Во-вторых, точное вычисление платежей может быть дорогим. Для больших систем
необходимы приближённые расчёты, ограничение числа кандидатов или кластерный
режим.

В-третьих, качество обучаемой политики зависит от репрезентативности
обучающих сценариев. Если реальные потоки задач и сетевые задержки сильно
отличаются от симуляционной среды, требуется дообучение.

В-четвёртых, кластеризация снижает накладные расходы, но может ухудшать
качество распределения, если сильные межкластерные связи оказываются
разорванными.

\section{Выводы по главе}

В главе предложен гибридный метод AMPLE-AMR для адаптивного распределения
вычислительных задач автономных мобильных роботов между периферийными
вычислительными узлами. Метод объединяет обучаемую многоагентную политику
выбора операционных режимов узлов и VCG-подобный аукционный слой
распределения задач.

Ключевым результатом является строгое разделение ролей: обучаемый контур
управляет режимами допуска, резервирования и обслуживания задач, а
аукционный контур распределяет задачи в рамках выбранных ресурсных
ограничений. Это позволяет не обучать ставки и оценки ценности задач,
сохраняя экономическую интерпретируемость аукционного слоя.

Также предложена масштабируемая версия C-AMPLE-AMR, основанная на
кластеризации графа периферийной инфраструктуры и
проведении локальных аукционов. Введён критерий перехода от глобального
распределения к кластерному режиму на основе сравнения накладных расходов и
оценки возможной потери общественного благосостояния.

Полученные результаты образуют центральный методический вклад диссертации.
В следующей главе описывается программный прототип, реализующий основные
компоненты AMPLE-AMR и C-AMPLE-AMR для последующей экспериментальной
валидации.
